\begin{textsample}{POS Dim 1 – se – Score 39.00 – se/se\_inf\_e\_ef\_3.txt}  \label{ex:f1_pos_113}
O saber escolar é produto e produtor de valores.
Compete à escola, portanto, subordinar, criticamente, o seu projeto a \textbf{um} conjunto de princípios fundamental, à formação de \textbf{um} cidadão comprometido com o cuidado de si e do outro, comprometido com a democracia, com o princípio da diferença, com sua cidade, com o Brasil, com o desenvolvimento social, inclusivo e sustentável dessas \textbf{territorialidades}.
A opção por determinado programa de ensino não traz automaticamente esse ou aquele conjunto de valores. \textbf{Uma} filiação incidental não assegura \textbf{que} a escola se torne, no âmbito de cada comunidade, \textbf{um} polo formador de virtudes.
Não é suficiente estudar o conceito de democracia para se tornar democrático, por exemplo.
Por isso, fazem -se necessário \textbf{que} aos conteúdos se somem \textbf{uma} estética, \textbf{uma} ética e \textbf{um} \textbf{método} - ou a mescla de alguns caminhos - todos sempre alinhados ao \textbf{que} se deseja ensinar.
Relações autoritárias ou demasiadamente \textbf{hierarquizadas} nunca serão ambiente favorável para aprender sobre democracia.
Se a escola \textbf{supõe}, ao elaborar seu currículo, quais são os saberes \textbf{que} objetiva produzir, quais são as competências socioemocionais e cognitivas \textbf{que} deseja ensinar impõe -se, ao seu coletivo, selecionar os \textbf{métodos} mais apropriados para tais propósitos.
Organizar espaços para o conhecer, para o fazer, para o conviver e para estar e ser resumem a grande \textbf{tarefa} da escola básica \textbf{contemporânea}.
Não se cumpre esse desafio de qualquer modo.
Menos ainda sob a égide da espontaneidade.
Quais são as metodologias mais adequadas ao cumprimento desse projeto escolar?
Pensar e formular estratégias nessas \textbf{direções} constituem as artes dos educadores.
Elas serão adequadas se \textbf{conseguirem} relacionar os conteúdos escolares com os contextos dos \textbf{discentes}, se se \textbf{basearem} numa \textbf{pedagogia} da pergunta, se comprometerem o estudante com valores atitudinais, se forem \textbf{dialógicas} e se oportunizarem a cada \textbf{discente} o sentimento da co-responsabilização e o exercício da autoaprendizagem.
Toda criança é capaz de aprender.
Todo adolescente é capaz de aprender, mesmo quando diagnosticado com algum nível de deficiência cognitiva.
Nesse último caso, as aprendizagens, obviamente, dar -se-ão, como em qualquer outra circunstância, nos limites estabelecidos pelo corpo.
Com ou sem deficiência serão as estratégias \textbf{metodológicas} \textbf{que} cumprirão papel fundamental.
A escola básica está dividida em níveis, etapas e modalidades de ensino.
Passando pela escolarização num fluxo adequado, o estudante e a estudante \textbf{atravessarão} esses tempos, em fases específicas do desenvolvimento socioemocional e psíquico-motor.
O currículo deve reconhecer os limites, as potencialidades e desafios comumente experimentados na infância, na adolescência e na juventude.
E deve ser suficientemente plástico para dar conta das adversidades e desencontros nesse fluxo reconhecendo, ainda, a historicidade das \textbf{categorias}, anteriormente mencionadas.
Quem é o \textbf{sujeito} da infância, da adolescência e da juventude dos dias \textbf{contemporâneos}?
Quem são eles, em cada \textbf{uma} das escolas sergipanas?
Quem é o estudante para quem se pensa essa rede de saberes?
É imperioso \textbf{que} cada coletivo escolar se ocupe dessas e de outras perguntas ao se colocar na \textbf{tarefa} de elaboração do seu currículo.
Nessa perspectiva, a escola deve reconhecer, pelo menos, o duplo tempo do estudante : o tempo singular em \textbf{que} ele produz sua aprendizagem e o tempo histórico-social nos quais ele vai se constituindo como sujeito.
A cibercultura, por exemplo, é \textbf{um} território \textbf{que} não se pode negligenciar.
Impõe -se \textbf{um} esforço de todos os educadores para \textbf{que} a escola fale também no e do tempo das linguagens e do universo virtual com os quais as mais recentes gerações vêm se constituindo.
As crianças e adolescentes \textbf{vivem} \textbf{uma} experiência singular com o fenômeno das linguagens.
Nunca foram tantas atingindo e produzindo, simultaneamente, cada \textbf{sujeito}.
Ao projeto escolar, portanto, compete a construção de espaços de estudos \textbf{que} capacitem os \textbf{discentes} para a leitura e produção de diferentes linguagens.
Essas questões adquirem absoluto sentido quando se reconhece, dentre outras características, \textbf{que} o Brasil, na sua continentalidade, é \textbf{uma} miríade de “ nações ”, de “ paisagens socioculturais ”.
A escola e seu currículo devem legitimar e promover essa diferença.
O campo e a cidade, o urbano e o rural produzem subjetividades e geografias distintas.
A história de vários povos tradicionais pulsa, com mais ênfase, em diversos pontos do território brasileiro.
O currículo deve contemplá -las.
As culturas camponesas, extrativistas, quilombolas, \textbf{ribeirinhas} e dos povos das florestas, por exemplo, devem reverberar -se nesse projeto escolar.
Seja na instituição \textbf{que} se constitua no âmbito de qualquer dessas comunidades, “ \textbf{ancorando} -se em suas temporalidades e saberes ”, seja naquela \textbf{que}, estando em qualquer lugar, oportunize \textbf{um} ensino \textbf{que} contemple essas diferenças.
Construir \textbf{um} projeto educacional com todos esses \textbf{compromissos} \textbf{exige} daqueles \textbf{que} o fazem a observação e análise de todos os processos.
Daqueles \textbf{que} acontecem no âmbito de cada sala, na relação professor/aluno, destes com os objetos de conhecimento, mas também de todos aqueles \textbf{que} devem garantir o funcionamento satisfatório da escola.
Avaliar torna -se \textbf{uma} \textbf{tarefa} permanente.
Avaliar com o propósito de garantir \textbf{que} todos os esforços convirjam para produzir \textbf{sujeitos} compromissados com os princípios e valores \textbf{basilares} desse projeto escolar.
Numa escola \textbf{que} reconhece e promove a diferença, os modos de mensurar o desenvolvimento dos processos pedagógicos, por exemplo, devem ser diversos e devem adequar -se às competências socioemocionais e cognitivas \textbf{que} se deseja ensinar.
Qualquer prática avaliativa deve constituir -se, antes de tudo, como mais \textbf{um} espaço para a aprendizagem.
Avalia -se, na escola democrática e inclusive, com vistas à qualificação e reelaboração, quando necessário, dos caminhos, das rotinas e, principalmente, com o propósito de acompanhar o desenvolvimento progressivo da aprendizagem.
Sob esse projeto escolar, não se deve avaliar para classificar, \textbf{hierarquizar} e reter.
Se todos são capazes de aprender, compete às práticas avaliativas apontar para os percursos mais apropriados ao \textbf{sujeito} \textbf{que} aprende e aos contextos em \textbf{que} se produzem as aprendizagens.
Diante desse panorama, impõe -se \textbf{um} chamado ao educador de cada escola básica no estado de Sergipe e, na mesma medida, convocam -se todas as instituições, a sociedade como \textbf{um} todo, com vistas à execução de \textbf{um} projeto curricular em cada escola \textbf{que} deverá ser construído numa perspectiva \textbf{dialógica} ; considerando o \textbf{que} dispõem os documentos \textbf{norteadores} da educação nacional, especialmente suas Diretrizes Curriculares e a BNCC ; tendo como foco a aprendizagem do aluno, garantida como direito humano de todos e de cada \textbf{um}, respeitando, assim, os tempos das aprendizagens e as \textbf{singularidades} desses processos em cada unidade de ensino.
Sabe -se do tamanho do desafio.
Impõe -se a mobilização dos \textbf{atores} na (re)invenção de \textbf{um} currículo desenhado com e para a unidade, cujo produto seja \textbf{uma} escola democrática \textbf{que} permita ao sujeito estudante compreender e desenvolver conceitos-ferramentas para intervenção e produção de sua existência material e simbólica.
Esse currículo, entretanto, deve estar compromissado com a promoção da diferença, com o exercício do diálogo, com a construção de cidades socialmente desenvolvidas e, portanto, inclusivas.
No \textbf{bojo} dos saberes socialmente produzido e \textbf{historicamente} acumulado deve a escola valorizar e problematizar, à luz da legislação educacional, as diversas linguagens artísticas, a diversidade religiosa e as outras tantas formas de experimentação do sobrenatural ; deve reconhecer a dimensão histórico-social das construções de gênero e de orientação sexual, do mesmo modo, as implicações dos processos culturais de “ normatizações ”.
Deve valorizar a diversidade cultural e étnica \textbf{que} se apresentam na vida social da cidade, do estado e do país ; deve \textbf{educar} para o exercício da fruição ; para posicionar -se diante do outro com desenvoltura argumentativa, \textbf{baseada} em fatos, dados e em informações confiáveis ; deve prepará -los para o cuidado de si e do outro e deve reconhecer esse conjunto de coisas como \textbf{um} direito humano.
Os modos de estar no mundo \textbf{demarcam} \textbf{uma} estética, \textbf{uma} ética, \textbf{uma} posição política, mesmo quando não se tem consciência disso.
Não compete ao projeto escolar, em temas sensíveis como \textbf{sexualidade}, gênero e religiosidade defender esse ou aquele caminho.
Aos educadores, fica a \textbf{tarefa} de organizar e coordenar ambientes de aprendizagens \textbf{que} oportunizem às crianças e aos adolescentes o desenvolvimento da autoestima, o respeito por si e pelo outro, independente das travessias subjetivas \textbf{que} todo cidadão tem o direito de realizar, devendo, inclusive, intervir dentro dos parâmetros e dos propósitos educacionais, para \textbf{que} a caminhada de cada \textbf{um} não coloque mediante risco a \textbf{dignidade} de si e dos \textbf{que}, ao \textbf{lado}, fazem também sua jornada.
O movimento é \textbf{uma} condição do nosso tempo. (Re)significar a escola brasileira se impõe.
Torná -la efetivamente democrática, reorganizando -a para a educação do \textbf{século} XXI é a \textbf{tarefa} de todos.
E essa \textbf{tarefa} somente se cumprirá, em cada escola, por meio da elaboração e execução de \textbf{um} currículo.
É \textbf{urgente} oferecer \textbf{um} ensino das “ ciências ”, das linguagens e da matemática submerso em valores e alinhado com \textbf{uma} efetiva preparação para o exercício crítico da cidadania.
E cidadania inclui dentre tantas outras coisas, a preparação para o mundo do trabalho.
O trabalho produziu o \textbf{sujeito} humano e continua a nos reinventar.
Não adianta, tão somente, garantir o acesso.
É fundamental \textbf{que}, nessa escola, as crianças e jovens aprendam, \textbf{que} sejam \textbf{aprendizes} partícipes, \textbf{que} dividam com os professores e gestores, mas sob a tutela destes, o protagonismo de \textbf{uma} escola cidadã.
Para dar conta desse desafio, além das condições objetivas, fundamentais ao funcionamento da escola, ao trabalho diário dos seus educadores, é condição \textbf{basilar} o \textbf{compromisso} ético dos gestores da rede e de todos aqueles \textbf{que} fazem a escola cidadã.
É \textbf{imprescindível} \textbf{que} \textbf{educar} crianças e adolescentes seja \textbf{um} projeto de cada cidade, da sociedade.
O desafio de aprender para ensinar e aprender se apresentará cotidianamente.
Por isso, a formação permanente no âmbito da escola, submersa em seu coletivo, se impõe e é ela, sobretudo, \textbf{que} assegurará o vínculo entre o projeto, suas metas e objetivos, e concederá a cada educador \textbf{um} papel sem igual nessa \textbf{tarefa} de construção de tantos projetos de vida, contribuindo, dessa forma, para a reinvenção de cada urbe, numa perspectiva de cidade sustentável, inteligente, humana e criativa.

% matched lemmas: ancorar, aprendiz, ator, atravessar, basear, basilar, bojo, categoria, compromisso, conseguir, contemporâneo, demarcar, dialógico, dignidade, direção, discente, educar, exigir, hierarquizar, historicamente, imprescindível, lado, metodológico, método, norteador, pedagogia, que, ribeirinho, sexualidade, singularidade, sujeito, supor, século, tarefa, territorialidade, um, uma, urgente, viver
\end{textsample}
